\chapter*{Lampiran}
\addcontentsline{toc}{chapter}{Lampiran}

\section*{A. Rubrik Penilaian Tugas Praktikum (Ringkas)}

\par\addvspace{\ModulMediaskip}
\begin{longtable}{|>{\raggedright\arraybackslash}p{3.2cm}|>{\raggedright\arraybackslash}p{4.2cm}|>{\raggedright\arraybackslash}p{6.6cm}|}
\caption{Rubrik penilaian tugas praktikum (ringkas)}\label{tab:lamp-rubrik}\\
\hline
\textbf{Aspek} & \textbf{Bobot indikatif} & \textbf{Indikator} \\
\hline
\endfirsthead
\hline
\textbf{Aspek} & \textbf{Bobot indikatif} & \textbf{Indikator} \\
\hline
\endhead
\hline
\endfoot

Kebenaran program & 40\% & Kompilasi tanpa error fatal; keluaran sesuai spesifikasi; penanganan kasus uji \\
\hline
Struktur kode & 25\% & Fungsi terpisah wajar; modularitas \texttt{.c}/\texttt{.h} bila diminta; tidak ada duplikasi berlebihan \\
\hline
Gaya dan dokumentasi & 15\% & Penamaan jelas; \texttt{-Wall -Wextra} bebas peringatan kritis; \texttt{README} singkat jika diminta \\
\hline
Laporan / log & 20\% & Langkah jelas; cuplikan perintah \texttt{gcc}/\texttt{make}; tangkapan layar output bila diminta \\
\hline
\end{longtable}
\addvspace{\ModulMediaskip}

\section*{B. Daftar Referensi dan Sumber Belajar Terbuka}

Sumber berikut mendukung penyusunan modul ini (selaras \texttt{sumber\_materi.md} dan \texttt{silabus\_praktikum\_c.tex}):

\begin{enumerate}[leftmargin=*]
  \item UNESA. \textit{RPS Praktikum Pemrograman Dasar} (PDF). \url{https://sindig.unesa.ac.id/rps-pdf/d4-manajemen-informatika/prak-pemrograman-dasar.pdf}
  \item UNS. \textit{RPS Pemrograman Dasar} EE0107-19 (PDF). \url{https://elektro.ft.uns.ac.id/wp-content/uploads/2020/08/RPS_EE0107-19_Pemrograman-Dasar.pdf}
  \item UNWAHAS. \textit{RPS Praktikum Algoritma dan Pemrograman} JTI2403. \url{https://simpel.teknik.unwahas.ac.id/rps/191}
  \item cppreference.com. \textit{C language and library reference}. \url{https://en.cppreference.com/w/c}
  \item Learn-C.org. \textit{Interactive C tutorial}. \url{https://www.learn-c.org/}
  \item The Little Book of C. \url{https://little-book-of.github.io/c/books/en-US/book.html}
  \item Petani Kode. \textit{Tutorial Pemrograman C}. \url{https://petanikode.com/tutorial/c}
  \item W3Schools. \textit{C Tutorial}. \url{https://www.w3schools.com/c/index.php}
  \item Programiz. \textit{C Programming}. \url{https://www.programiz.com/c-programming}
  \item GeeksforGeeks. \textit{C Programming language}. \url{https://www.geeksforgeeks.org/c-programming-language/}
  \item Harvard CS50x. \textit{Week 1 --- C}. \url{https://cs50.harvard.edu/x/2026/weeks/1/}
  \item Beej. \textit{Beej's Guide to C Programming}. \url{https://beej.us/guide/bgc/html/index.html}
  \item GNU Project. \textit{The GNU C Reference Manual}. \url{https://www.gnu.org/software/gnu-c-manual/gnu-c-manual.html}
  \item FSF. \textit{Using the GNU Compiler Collection (GCC)}. \url{https://gcc.gnu.org/onlinedocs/gcc/index.html}
  \item Gustedt, J. \textit{Modern C}. \url{https://gustedt.gitlabpages.inria.fr/modern-c/}
\end{enumerate}

\section*{C. Evaluasi Akhir dan Integrasi Kompetensi}

Dosen dapat menggabungkan: (1) portofolio tugas mingguan Bab 1--12; (2) UTS praktikum (Minggu 14) memetakan CPMK-1--4; (3) mini-proyek multi-berkas dengan \texttt{Makefile} (Minggu 15); (4) UAS praktikum (Minggu 16) berfokus CPMK-3--6. Rubrik komprehensif menautkan indikator ke Sub-CPMK pada \texttt{silabus\_praktikum\_c.tex}.

\section*{D. Komponen Penilaian Semester}
\noindent\textit{Selaras dokumen \path{silabus_praktikum_c.tex}.}
\vspace{0.35em}

\begin{modultable}[]{Komponen penilaian semester (bobot indikatif)}
\begin{tabular}{|l|c|>{\raggedright\arraybackslash}p{7.2cm}|}
\hline
\textbf{Komponen} & \textbf{Bobot} & \textbf{Catatan / pemetaan} \\
\hline
Tugas praktikum mingguan & 28\% & CPMK-1 s.d. CPMK-6 sesuai minggu \\
\hline
Laporan / log praktikum & 17\% & Bukti Sub-CPMK; format sesuai dosen \\
\hline
Kinerja dan rubrik sesi & 10\% & Sikap, K3, partisipasi (CPL-4); observasi lab \\
\hline
Ujian praktikum tengah & 20\% & CPMK-1--4; Minggu 14 \\
\hline
Ujian praktikum akhir / proyek & 25\% & CPMK-3--6; Minggu 15--16 \\
\hline
\textbf{Jumlah} & \textbf{100\%} & \\
\hline
\end{tabular}
\end{modultable}

\section*{E. Lembar Observasi Asisten (contoh)}

Nama mahasiswa / NIM: \ldots\quad Tanggal: \ldots\quad Bab / Sub-CPMK: \ldots

\begin{itemize}[leftmargin=*, label=$\square$]
  \item K3 dan tata tertib laboratorium dipatuhi (CPL-4)
  \item Prosedur kompilasi/\texttt{make} diikuti tanpa intervensi berlebihan
  \item Menunjukkan pemahaman konsep (bukan hanya menyalin kode)
  \item Artefak \texttt{.c}/\texttt{.h} tersimpan dan dinamai sesuai konvensi
\end{itemize}

\noindent Komentar: \hrulefill

\section*{F. Template Penamaan dan Pengumpulan}

\begin{itemize}[leftmargin=*]
  \item Berkas sumber: \texttt{NIM\_TIF1xx\_M\textit{nn}\_tugas.c} (ganti \texttt{TIF1xx} dengan kode MK resmi prodi)
  \item Proyek akhir: folder \texttt{NIM\_TIF1xx\_proyek/} berisi \texttt{README.md}, \texttt{Makefile} (jika diminta), \texttt{src/}, \texttt{include/}
  \item Laporan: PDF dengan identitas dan ringkasan perintah \texttt{gcc -Wall -Wextra -std=c17}
\end{itemize}

\section*{G. Rangkuman Sub-CPMK dan Refleksi}

\begin{itemize}[leftmargin=*]
  \item \textbf{Rangkuman:} setelah Bab 12, mahasiswa telah menempuh Sub-CPMK 1.1 s.d. 6.3 sesuai RPS.
  \item \textbf{Refleksi (panduan):} tuliskan dua poin perkembangan terbesar terkait pointer/berkas; satu risiko yang Anda hindari di lab.
  \item \textbf{Umpan balik modul:} \ldots\ (angket singkat opsional prodi)
\end{itemize}

\section*{H. Glosarium Singkat}

\begin{description}[style=nextline, widest=Sub-CPMK:, itemsep=0.25em]
  \item[CPMK:] Capaian pembelajaran mata kuliah.
  \item[CPL:] Capaian pembelajaran lulusan program studi.
  \item[GCC:] GNU Compiler Collection (\texttt{gcc}).
  \item[GDB:] GNU Debugger.
  \item[OBE:] Pembelajaran berbasis capaian.
  \item[RPS:] Rencana pembelajaran semester.
  \item[Sub-CPMK:] Indikator operasional yang menurunkan CPMK.
  \item[Toolchain:] rangkaian kompilator, penaut, dan perangkat bantu (\texttt{make}, \texttt{gdb}).
\end{description}
